\documentclass[11pt]{article}
\usepackage[margin=1in]{geometry}
\usepackage[T1]{fontenc}
\usepackage[utf8]{inputenc}
\usepackage{lmodern,microtype,amsmath,amssymb,amsthm,mathtools}
\usepackage[colorlinks=true,linkcolor=blue,citecolor=blue,urlcolor=blue]{hyperref}
\hypersetup{pdftitle={Exponential online mistakes for isotropic Bayesian threshold networks at all sufficiently large widths},pdfauthor={ChatGPT-6 Astra; Kaarel Hanni; Dmitry Vaintrob}}
\newtheorem{theorem}{Theorem}
\newtheorem{lemma}[theorem]{Lemma}
\newtheorem{proposition}[theorem]{Proposition}
\newtheorem{corollary}[theorem]{Corollary}
\newtheorem{remark}[theorem]{Remark}
\newcommand{\R}{\mathbb R}
\newcommand{\E}{\mathbb E}
\newcommand{\Pp}{\mathbb P}
\newcommand{\sgn}{\operatorname{sign}}
\newcommand{\cH}{\mathcal H}
\title{Exponential online mistakes for isotropic Bayesian\protect\\threshold networks at all sufficiently large widths}
\author{ChatGPT-6 Astra \and Kaarel H\"anni \and Dmitry Vaintrob\thanks{Authors are listed in alphabetical order.}}
\date{6 November 2026}
\begin{document}
\maketitle

\begin{abstract}
For the fixed discontinuous sign activation, we construct an exactly
isotropic Gaussian Bayesian network with exponentially many online
mistakes at every sufficiently large hidden width.
The model retains one hidden layer, a constant input coordinate, the full
Boolean cube, global output margin greater than $1/2$, and exact labels.
For each $k\ge2$ the target has a Boolean circuit of size $\Theta(k)$ and
a width-one teacher with every weight at most one, while one fixed input
permutation forces $2^k-2$ designated mistakes after an exponential
calibration prefix. The result holds under every fixed strong-prior
exponent $\gamma>1$, with any fixed positive variance prefactor, at all
sufficiently large widths. It also holds for $\gamma=1$ with a sufficiently
small fixed prefactor. A finite-width Gaussian-tail inequality controls
every nonlinear competitor and reduces the posterior to strictly ordered
threshold-function cone probabilities. An explicit positive-helper
calibration and radial integration produce that order without changing
the isotropic covariance. The same target and prior have polynomial IID
sample complexity. The cold-posterior width bound is
$\operatorname{poly}(k)2^k$; the direct fixed-$\gamma$ theorem has an
additional width requirement stated explicitly. This does not establish
the corresponding result for ordinary tanh, ReLU, GELU, or every width
above the teacher width.
\end{abstract}

\section{Model and theorem}

Let $X=\{-1,1\}^d$, write $\bar x=(1,x)$, and use the fixed activation
$\sgn(t)=1$ for $t>0$ and $-1$ for $t<0$. Its value at zero may be fixed
arbitrarily: no Gaussian first-layer score is zero on the finite cube
with positive probability. The learner has one hidden layer,
\begin{equation}\label{eq:network}
 F_{W,v}(x)=\sum_{a=1}^N v_a\sgn(W_a\cdot\bar x),
 \qquad
 W_{ai}\sim N(0,N\sigma^2/d),\quad
 v_a\sim N(0,\sigma^2/N),
\end{equation}
with all variables independent, including the bias coordinate $i=0$.
There is no coordinate-dependent covariance. For a Boolean function $g$
let
\begin{equation}\label{eq:cell}
 U_g=\{g(x)F_{W,v}(x)>1/2\text{ for every }x\in X\}.
\end{equation}
The posterior conditions the prior on the union of these strict
global-margin cells and on the exact observed labels. Thus it is a
posterior over Boolean functions on the whole cube, not a restriction
to threshold subnetworks.

By a posterior error probability at an input we mean the posterior mass
of functions assigning the wrong label. Posterior-label majority makes
a deterministic mistake whenever this probability exceeds $1/2$.
The stronger threshold $2/3$ will also be covered. A fresh posterior
function sampled for prediction errs with that posterior probability.

For a fixed $\gamma\ge1$ and $\tau>0$, the strong-prior scaling is
\begin{equation}\label{eq:gamma}
 \sigma_N^2=\tau^2N^{-\gamma},\qquad
 \operatorname{Var}(W_{ai})=\frac{\tau^2N^{1-\gamma}}d,
 \qquad \operatorname{Var}(v_a)=\tau^2N^{-\gamma-1}.
\end{equation}
The unit-prefactor version sets $\tau=1$.

\begin{theorem}\label{thm:main}
For every $k\ge2$, let $d=167k+1617$. There is an explicitly specified
threshold function $f_k:X\to\{-1,1\}$ and a fixed permutation of $X$
with the following properties.
\begin{enumerate}
\item The target has a bounded-fan-in Boolean circuit of size
$\Theta(d)=\Theta(k)$ and a width-one sign-network teacher with every
weight at most one in absolute value and output margin one.
\item The permutation begins with $2^{d-1}+2$ calibration examples,
followed by $2^k-2$ distinct designated inputs whose true labels are
negative. For every fixed $\gamma>1$, $\tau>0$, and $\eta>0$, there is
$N_0$ such that at every $N\ge N_0$ under \eqref{eq:gamma}, the posterior
error probability at every designated round exceeds $1-\eta$.
\item Consequently posterior-label majority makes all $2^k-2$
designated mistakes, including under the strict $2/3$ criterion, at all
sufficiently large widths. The probability that posterior sampling makes
all these mistakes can also be made greater than $1-\eta$ by increasing
the same width threshold.
\item The conclusion also holds at $\gamma=1$ for every sufficiently
small fixed $\tau>0$. It is not asserted at $\gamma=1,\tau=1$.
\end{enumerate}
The target, input order, and first-layer isotropy do not depend on $N$,
$\gamma$, or $\tau$. The threshold $N_0$ may depend on $k$, $\gamma$,
$\tau$, and the required error allowance.
\end{theorem}

This answers an all-sufficiently-large-width question for a classical
discontinuous activation. It leaves the exactly isotropic tanh question
open. The small-anisotropy tanh result of~\cite{anisotropic} preserves
ordinary tanh but changes the covariance; the present result preserves
the covariance and changes the activation. The exponential calibration
prefix and the large width threshold are substantive features of the
construction, not efficient-data or practical-width guarantees.

\section{A finite-width posterior bound}

Let $\cH_d$ be the affine threshold functions on $X$ and let
$\pi(g)$ be the probability that one isotropic Gaussian hidden neuron
has truth table $g$. Set $\pi(g)=0$ for $g\notin\cH_d$.
For $g\in\cH_d$, its strict representation cone is nonempty and open,
so $\pi(g)>0$. Gaussian symmetry gives $\pi(-g)=\pi(g)$.
Positive rescaling of every first-layer coefficient does not change
the hidden truth table, so these probabilities are independent of the
common Gaussian variance and of $N$.

Write $Q(a)=\Pp(Z>a)$ for a standard Gaussian $Z$. Set
\begin{equation}\label{eq:tailconstants}
 T=2^d,\quad c_d=1-(1-2/T)^2,\quad p_g=2\pi(g),\quad
 q_\sigma=Q(1/(2\sigma)),\quad
 \delta_{N,\sigma}=e^{-c_d/(8N\sigma^2)}.
\end{equation}

\begin{lemma}[All Boolean competitors]\label{lem:posterior}
For every Boolean $g$, every $N\ge1$, and every $\sigma>0$,
\begin{equation}\label{eq:massbounds}
 q_\sigma p_g^N\ \le\ \Pp(U_g)\ \le\
 q_\sigma[p_g+(1-p_g)\delta_{N,\sigma}]^N.
\end{equation}
\end{lemma}
\begin{proof}
Condition on the hidden truth tables $h_1,\ldots,h_N$. Let $t$ count
those not equal to either $g$ or $-g$, and put
\[
 m_a=\frac1T\sum_{x\in X}g(x)h_a(x).
\]
Aligned columns have $|m_a|=1$. Every other column both agrees and
disagrees with $g$, hence $|m_a|\le1-2/T$. Thus
$\|m\|^2\le N-c_dt$. If $U_g$ occurs, averaging its strict inequalities
gives $v\cdot m>1/2$. This scalar Gaussian has variance
$\sigma^2\|m\|^2/N$, so
\[
 \Pp(U_g\mid h_1,\ldots,h_N)
 \le Q\!\left(\frac1{2\sigma\sqrt{1-c_dt/N}}\right)
 \le q_\sigma e^{-c_dt/(8N\sigma^2)}.
\]
For the last inequality use $(1-r)^{-1}-1\ge r$ and
$Q(b)/Q(a)\le e^{-(b^2-a^2)/2}$ for $b\ge a\ge0$.
The latter follows by shifting the Gaussian-tail integral by $b-a$.
If all columns are aligned, $U_g$ is exactly one scalar tail of
probability $q_\sigma$. Alignment occurs independently with probability
$p_g$ per neuron. Restricting to $t=0$ proves the lower bound; averaging
the upper bound over its binomial distribution proves the upper bound.
\end{proof}

\begin{corollary}[Cold posterior]\label{cor:cold}
At each fixed width $N$, after any exact-label history consistent with
some threshold function, the posterior as $\sigma\downarrow0$ gives
mass proportional to $\pi(g)^N$ on the consistent threshold functions
and zero mass on the other Boolean functions.
\end{corollary}
\begin{proof}
Divide \eqref{eq:massbounds} by $q_\sigma$ and let $\sigma\downarrow0$.
There are finitely many Boolean functions and at least one positive
limiting mass. The common factor $2^N$ cancels on normalization.
\end{proof}

\begin{proposition}[Direct strong-prior limit]\label{prop:gamma}
Fix any finite collection of exact-label histories. If each has a unique
largest-$\pi$ consistent threshold function, then under
\eqref{eq:gamma}, for every fixed $\gamma>1$ and $\tau>0$, the posterior
mass of that function tends to one as $N\to\infty$, simultaneously
at all the histories. The same holds for $\gamma=1$ and all sufficiently
small fixed $\tau>0$.
\end{proposition}
\begin{proof}
For $\gamma>1$,
\[
 \delta_N=\exp[-c_dN^{\gamma-1}/(8\tau^2)],\qquad
 N\delta_N\longrightarrow0.
\]
For every threshold function, \eqref{eq:massbounds} gives
$\Pp(U_g)/(q_{\sigma_N}p_g^N)\to1$.
For a nonthreshold function it gives
$\Pp(U_g)\le q_{\sigma_N}\delta_N^N$.
Dividing by the lower bound for the largest consistent threshold
function proves the claim after summing the finitely many other cells.

For $\gamma=1$, $\delta_N=\delta_\tau=e^{-c_d/(8\tau^2)}$ is constant.
At one history let $p_*$ be the largest aligned probability and $p_2$
the largest among the other consistent threshold functions, taking
$p_2=0$ if none exist. Include the nonthreshold competitors by the same
zero value. Choosing
\[
 \delta_\tau<\frac{p_*-p_2}{1-p_2}
\]
makes every competing upper-bound base in \eqref{eq:massbounds}
strictly smaller than the lower-bound base $p_*$. The finite collection
of histories permits one common positive upper bound on $\tau$.
\end{proof}

\section{An explicit full-cube calibration}

The next construction supplies strictly ordered Gaussian cone
probabilities. No parameter is observed or conditioned directly:
calibration consists only of exact labels of one threshold function.

\subsection{Four registers and positive helper pairs}

Fix $k\ge2$, $p\ge6$, and $L=p+k+2$. Consider four logical $L$-bit
registers $A,B,C,D$, with levels $0,\ldots,L-1$. The actual main inputs
are the $k$ bits of $A$ at levels $p,\ldots,p+k-1$.
Every other logical bit is represented by a pair of actual helper inputs.
There are
\begin{equation}\label{eq:dimensions}
 P=4L-k\text{ pairs},\qquad H=2P=8L-2k\text{ helpers}.
\end{equation}
The signed logical coefficient is $+2^j$ at level $j$ in $A,B$ and
$-2^j$ in $C,D$. Denote the signed coefficient for helper pair $l$ by
$r_l^*$. The canonical main coefficients are $w_i^*=2^{p+i}$, and put
\begin{align}\label{eq:canonical}
 R_*&=(2^k-1)2^p,&
 S_*&=\sum_l|r_l^*|=4(2^L-1)-R_*,\\
 H_*&=R_*+S_*+1=2^{L+2}-3,& u_*&=1/4.\nonumber
\end{align}
The two actual helper coefficients of pair $l$ are
$H_*+r_l^*/2$ and $H_*-r_l^*/2$.
Let $f_0(x,y)$ be the sign of this affine score on the entire cube of
main and helper inputs.

All scores are nonzero: each pair contributes either $\pm2H_*$ or
$\pm r_l^*$, so the nonbias score is integral. If the total helper-bit
sum $m=\sum y_l$ is nonzero, $|m|\ge2$, and the label is $\sgn(m)$,
because $H_*m$ dominates the other terms. If every pair is individually
balanced, its inputs are $(s,-s)$ and its contribution is $r_l^*s$.
On this subcube the truth table is exactly
\begin{equation}\label{eq:addition}
 A+B\ge C+D,
\end{equation}
where register values are interpreted in $0/1$ notation. Equality is
positive because $u_*>0$.

Now consider \emph{any} strict affine representation of the full truth
table $f_0$. Write its main coefficients as $w_i$, helper coefficients
as $h_l$, and bias as $u$. Every assertion below applies to this whole
representation cone.

\subsection{Arithmetic rigidity}

\begin{lemma}\label{lem:pinning}
Let $a=w_0$ and $R=\sum_iw_i$. Then $a>0$, all $w_i>0$, and every binary
subset sum $A_j=\sum_i\operatorname{bit}_i(j)w_i$, $0\le j\le2^k-1$,
satisfies
\begin{equation}\label{eq:binary}
 \left|A_j/a-j\right|<\epsilon,
 \qquad \epsilon=12\,2^{k-p}.
\end{equation}
Every within-pair coefficient difference has absolute value at most
$10R$, and $|u|\le4LR$.
\end{lemma}
\begin{proof}
On the individually pair-balanced subcube, use $a_j,b_j,c_j,d_j$ for
the positive register coefficients, writing negative signs for $C,D$.
They are strictly positive because every logical coordinate in
\eqref{eq:addition} is essential and has its indicated monotonicity.
In particular $w_i=a_{p+i}$. After switching to $0/1$ coordinates and
dividing the score by two, it is
\[
 \beta+\sum_j a_jA_j+\sum_j b_jB_j-\sum_jc_jC_j-\sum_jd_jD_j,
 \quad
 \beta=\frac{u-\sum a_j-\sum b_j+\sum c_j+\sum d_j}{2}.
\]
The all-zero input and the inputs $C=1$ or $D=1$ imply
$0<\beta<\min(c_0,d_0)$. Put $u_0=\max(c_0,d_0)$.
Compare the equality $A=2^j,C=2^j$ first with $D=0$ and then with $D=1$.
The same comparison works with $B$ in place of $A$, giving
\begin{equation}\label{eq:copies}
 -\beta<a_j-c_j<d_0-\beta,
 \qquad -\beta<b_j-c_j<d_0-\beta.
\end{equation}
Similarly $a_j-d_j,b_j-d_j\in(-\beta,c_0-\beta)$.
All coefficients at a given level differ from $a_j$ by less than $u_0$,
and $|b_j-a_j|<d_0$.
For $j<L-1$ compare $A=B=2^j,C=2^{j+1}$ first with $D=0$ and then
with $D=1$. These labels imply
\begin{equation}\label{eq:recurrence}
 \beta-d_0<c_{j+1}-a_j-b_j<\beta,
 \qquad |a_{j+1}-2a_j|<3u_0.
\end{equation}

For $j\ge1$, compare $A=2^j-1$ with $C=2^j$ and with $C=2^j-1$.
The first label is negative and the second positive, so
$c_j>\beta+\sum_{i<j}a_i>\sum_{i<j}c_i$.
The same argument gives $d_j>\sum_{i<j}d_i$. Hence
\[
 a:=a_p>(2^{p-1}-1)u_0\ge2^{p-2}u_0.
\]
Iterating \eqref{eq:recurrence} now gives
$|a_{p+i}-2^ia|<3u_0(2^i-1)$.
The sum of the $k$ coordinate errors is less than $3u_02^k$, which
proves \eqref{eq:binary}.

For the coarse coefficient bounds, $u_0<a/16$ since $p\ge6$.
Above level $p$, $a_j<(19/16)2^{j-p}a$, while
$R>(13/16)(2^k-1)a$. Thus even at $j=L-1=p+k+1$,
$a_j/R<(38/13)2^k/(2^k-1)\le76/13<6$.
The other register coefficients differ from $a_j$ by less than $u_0$.
Below level $p$, $c_j<c_p<a+u_0$ and
$a_j<c_j+u_0<a+2u_0$; all other coefficients there are less than $2R$.
Thus $10R$ bounds every register coefficient and every absolute
within-pair difference. Finally \eqref{eq:copies} and its counterpart
for $b_j-d_j$ give
\[
 |u|=\left|2\beta+\sum_j(a_j-c_j)+\sum_j(b_j-d_j)\right|
 <(2L+2)u_0\le4LR.
\]
\end{proof}

\subsection{All helpers are positive and have bounded deviations}

\begin{lemma}\label{lem:helpers}
Every helper coefficient satisfies $h_l>R$. Write
\[
 H_0=H^{-1}\sum_lh_l,\qquad \delta_l=h_l-H_0,\qquad\sum_l\delta_l=0.
\]
Then $|\delta_l|\le30LR$ for every $l$.
\end{lemma}
\begin{proof}
Let $s_l$ be the mean of the two coefficients in helper pair $l$.
Fix two distinct pairs $i,j$, set pair $i$ both positive and pair $j$
both negative, and set every main input positive. Balance all other
pairs individually, choosing first all their canonical logical
contributions positive and then all negative. Both assignments have
total helper-bit sum zero. Their remaining canonical absolute sum is
at least
\[
 S_*-2\,2^{L-1}=3\,2^L-R_*-4>R_*+1/4,
\]
so the two labels are respectively positive and negative.
Arithmetic rigidity fixes the signs of the actual pair differences.
If $S_{ij}$ is the sum of their absolute values outside pairs $i,j$,
these two strict inequalities give
\[
 |2(s_i-s_j)+u+R|<S_{ij}\le S:=\sum_l|h_{l,+}-h_{l,-}|\le40LR.
\]
Thus $|s_i-s_j|\le(22L+1)R$. Averaging over $j$ and adding half a
within-pair difference gives $|h_l-H_0|\le30LR$.

To prove the stronger positivity statement, fix one actual helper
coordinate $t$. Balance its pair with that coordinate negative, and
balance every other pair so that its canonical logical contribution is
negative. The total canonical helper contribution is at most
$-S_*+2\,2^{L-1}<-R_*-1/4$. Thus the example with all main inputs
positive has negative label. Flip only helper coordinate $t$ to positive
and put every main input negative. Its total helper-bit sum is now $+2$,
so its label is positive. In any fitting representation, the difference
of the second score and the first is $2h_t-2R>0$. Hence $h_t>R$.
\end{proof}

\begin{lemma}[Exact common-mode ray]\label{lem:ray}
Let $\theta=(w,\delta,u)$ with $\sum\delta_l=0$. The full strict
representation cone of $f_0$ consists exactly of
\begin{equation}\label{eq:ray}
 \theta\in\Omega_0,\qquad H_0>A(\theta),
\end{equation}
where $\Omega_0$ is an open homogeneous cone independent of $H_0$ and
\begin{equation}\label{eq:raybounds}
 R\le A(\theta)\le BR,\qquad
 F_0:=\sum_iw_i^2+\sum_l\delta_l^2+u^2/2\le B^2R^2,
 \qquad B=200L^2.
\end{equation}
The binary approximation \eqref{eq:binary} holds on $\Omega_0$.
\end{lemma}
\begin{proof}
Define $\Omega_0$ by the labels on \emph{all} inputs with total
helper-bit sum zero, including assignments whose individual pairs are
not balanced. Those constraints do not involve $H_0$. On an input with
$m=\sum y_l\ne0$, the prescribed label is $\sgn(m)$; its constraint is
$H_0>-(u+w\cdot x+\delta\cdot y)/m$. Therefore \eqref{eq:ray} holds with
\[
 A(\theta)=\max_{x,y:\,m\ne0}
       -\frac{u+w\cdot x+\delta\cdot y}{m}.
\]
This is a finite maximum, so every $\theta\in\Omega_0$ has a nonempty
ray. All pinning and deviation bounds used only total-sum-zero examples
and hold on $\Omega_0$. The helper positivity conclusion holds at every
point above the ray threshold, hence
$A\ge R-\min_l\delta_l\ge R$.
Conversely $|m|\ge2$, $H\le8L$, and the preceding bounds imply
\[
 A\le(4L+1+240L^2)R/2\le125L^2R.
\]
Finally $\sum_iw_i^2\le R^2$, and
$F_0\le[1+7200L^3+8L^2]R^2\le(200L^2)^2R^2$.
\end{proof}

\section{Ordering the Gaussian solid angles}
\label{vol:section}

This section gives the geometric step of the construction. All Gaussian
coordinates in this section have the same variance; by homogeneity we take
that variance to be one. The calculation concerns the probability of an
\emph{entire threshold function}, rather than the volume of a selected set of
its representations.

There are $k$ main bits $x$, $H\ge2$ helper bits $y$, and one selector bit $z$.
Write an affine score as
\[
 b+c z+\sum_{i=0}^{k-1}w_i x_i+\sum_{\ell=1}^H h_\ell y_\ell,
 \qquad u=b-c,\quad v=b+c.
\]
Thus $u$ and $v$ are independent Gaussian variables of variance two,
independent of the unit-variance coordinates $w_i,h_\ell$.
Let $M=2^k-1$, and write
\[
 R=\sum_iw_i,\qquad A_j=\sum_i\mathrm{bit}_i(j)w_i\quad(0\le j\le M).
\]
In particular $A_0=0$ and $A_M=R$.

Suppose a complete truth table on the branch $z=-1$ gives an open homogeneous
representation cone $\mathcal C$ in $(w,h,u)$. We use the following properties
of the calibration constructed in the preceding section. Every point of
$\mathcal C$ has $a=w_0>0$, satisfies
\begin{equation}\label{vol:binary}
 \left|\frac{A_j}{a}-j\right|<\epsilon\quad(0\le j\le M),
 \qquad h_\ell>R\quad(1\le\ell\le H),
\end{equation}
and can be parametrized by
\[
 h_\ell=H_0+\delta_\ell,\qquad \sum_\ell\delta_\ell=0,
 \qquad \theta=(w,\delta,u)\in\Omega_0,
 \qquad H_0>A(\theta).
\]
Here $\Omega_0$ is an open homogeneous cone, independent of $H_0$, and
$A$ is homogeneous of degree one. For a constant $B\ge1$,
\begin{equation}\label{vol:scaffold-bounds}
 R\le A(\theta)\le BR,\qquad
 F_0:=\sum_iw_i^2+\sum_\ell\delta_\ell^2+u^2/2\le B^2R^2.
\end{equation}
In particular, the condition on $H_0$ is exactly a lower ray, with no upper
restriction. All these statements concern every representation in the
calibration cone.

\subsection{The compatible functions are exactly the cutoff functions}

Let $q_j$ be the point on the branch $z=+1$ whose helper bits are all $+1$
and whose main bits encode the integer $j$. Add the labels $q_0\mapsto-1$
and $q_M\mapsto+1$ to the calibration. Assume $\epsilon<1/2$, so the values
$A_j$ increase strictly with $j$.
The negative label at $q_0$ gives $v+\sum h_\ell-R<0$.
If at least one helper bit is negative, the affine score is at most
\[
 v+\sum h_\ell-2\min_\ell h_\ell+R
 <2R-2\min_\ell h_\ell<0.
\]
Consequently all threshold functions compatible with the calibration are
exactly the $M$ functions $g_0,\ldots,g_{M-1}$: they agree with the calibrated
function on $z=-1$, are negative off the helper face on $z=+1$, and satisfy
\[
 g_j(q_i)=-1\ \ (i\le j),\qquad
 g_j(q_i)=+1\ \ (i>j).
\]
Their full representation cones are
\begin{equation}\label{vol:full-cell}
 (w,h,u)\in\mathcal C,\qquad
 -\sum h_\ell+R-2A_{j+1}<v<-\sum h_\ell+R-2A_j.
\end{equation}
Each cell has positive Gaussian measure: $\mathcal C$ is nonempty and open,
and every displayed interval has positive length. Denote its Gaussian
probability by $p_j$.

\subsection{Radial integration and the common helper direction}

Set $D=k+H+1$ and $s=(D+1)/2$. Normalize $a=w_0$ to one.
There are $D+1$ original affine coefficients, and the radial change of
variables has Jacobian $a^D$. Integrating $a>0$ in the Gaussian density gives
a common positive constant times the power $-s$ of the normalized Gaussian
quadratic form. The changes $(b,c)\mapsto(u,v)$ and
$h\mapsto(H_0,\delta)$ contribute only constant Jacobians.
For fixed normalized $\theta$ put
\begin{equation}\label{vol:psi}
 \psi_\theta(r)=\int_{A(\theta)}^\infty
 \left[F_0+HH_0^2+\frac{(HH_0+r)^2}{2}\right]^{-s}\,dH_0.
\end{equation}
Equation~\eqref{vol:full-cell}, with $-v=HH_0+r$, therefore gives
\begin{equation}\label{vol:mass-integral}
 p_j=K\int_{\{\theta\in\Omega_0:w_0=1\}}
 \int_{-R+2A_j}^{-R+2A_{j+1}}\psi_\theta(r)\,dr\,d\theta,
\end{equation}
where $K>0$ and the cross-section measure are the same for every $j$.
All inner integration intervals lie in $[-R,R]$.

Define constants
\begin{align*}
 C_1&=B^2+H(B+1)^2+\frac{[H(B+1)+1]^2}{2},\\
 T_0&=2\max\{1,C_1/H\},\\
 C_2&=B^2+HT_0^2+\frac{(HT_0+1)^2}{2},\\
 \ell&=\frac{s(H-1)}{2C_2(M+1)}.
\end{align*}
The constants are polynomial in $H,B$, apart from the explicit factor
$(M+1)^{-1}$ in $\ell$.

\begin{lemma}[A uniform slope after integrating the common direction]
\label{vol:slope}
For every normalized $\theta$ and every $r\in[-R,R]$, if $\epsilon\le1$,
\[
 -\frac{\psi_\theta'(r)}{\psi_\theta(r)}\ge\ell.
\]
\end{lemma}
\begin{proof}
Write $F=F_0+HH_0^2+(HH_0+r)^2/2$.
The part of the integral in~\eqref{vol:psi} with
$A\le H_0\le A+R$ is at least $R(C_1R^2)^{-s}$.
The part with $H_0>T_0R$ is at most
\[
 \int_{T_0R}^\infty(HH_0^2)^{-s}\,dH_0
 =\frac{H^{-s}(T_0R)^{1-2s}}{2s-1}.
\]
The ratio of this upper bound to the preceding lower bound is
\[
 \frac{(C_1/H)^sT_0^{1-2s}}{2s-1}\le\frac12.
\]
Indeed $s\ge1$, and substituting $T_0=2\max\{1,C_1/H\}$ proves the inequality
separately for $C_1/H\le1$ and $C_1/H\ge1$.
Thus at least half the probability measure with density proportional to
$F^{-s}$ lies in $H_0\le T_0R$.
Differentiation under the integral is justified by the integrable power
bounds at infinity and gives
\[
 -\frac{\psi_\theta'(r)}{\psi_\theta(r)}
 =s\,\mathbb E\left[\frac{HH_0+r}{F}\right].
\]
The integrand is positive everywhere, because $H_0>A\ge R$ and $H\ge2$.
On $H_0\le T_0R$, its numerator is at least $(H-1)R$ and its denominator
at most $C_2R^2$. Hence the last display is at least
$s(H-1)/(2C_2R)$. Finally~\eqref{vol:binary} gives
$R=A_M<M+\epsilon\le M+1$.
\end{proof}

\begin{proposition}[Quantitative ordering of whole function cells]
\label{vol:ratio}
If
\[
 \epsilon\le\frac18,\qquad \epsilon\le\frac{\ell}{32},
\]
then, for $0\le j<M-1$,
\[
 p_j\ge(1+\ell/4)p_{j+1}.
\]
\end{proposition}
\begin{proof}
For fixed $\theta$, let $l_j=2(A_{j+1}-A_j)$.
By~\eqref{vol:binary}, $2-4\epsilon<l_j<2+4\epsilon$;
in particular $3/2\le l_j\le5/2$ and $l_{j+1}-l_j\le8\epsilon$.
Let $t$ be the common boundary of the two adjacent integration intervals.
Lemma~\ref{vol:slope} implies
\[
 \int_{t-l_j}^{t}\psi_\theta(r)\,dr
 \ge\psi_\theta(t)(l_j+\ell l_j^2/2),
 \qquad
 \int_t^{t+l_{j+1}}\psi_\theta(r)\,dr
 \le\psi_\theta(t)l_{j+1}.
\]
The difference of the first mass and the second is at least
\[
 \psi_\theta(t)(\ell l_j^2/2-8\epsilon)
 \ge(7\ell/8)\psi_\theta(t),
\]
whereas the second mass is at most $(5/2)\psi_\theta(t)$.
Thus the first is at least $(1+7\ell/20)$ times the second, which is stronger
than the asserted factor $1+\ell/4$. Integrate this inequality
in~\eqref{vol:mass-integral}.
\end{proof}

\subsection{A quantitative width bound in the cold posterior}

Take the target to be $g_{M-1}$ and, after calibration, reveal
$q_1,q_2,\ldots,q_{M-1}$, each with label $-1$.
There are $T=M-1=2^k-2$ designated rounds.
Suppose the cold posterior on surviving threshold functions is proportional
to $p_j^N$, as established by the posterior calculation for sign networks.
Immediately before $q_i$ is revealed, the survivors are
$g_{i-1},\ldots,g_{M-1}$. Exactly $g_{i-1}$ predicts the wrong label $+1$.
Writing $\rho=1+\ell/4$, Proposition~\ref{vol:ratio} and a geometric-series
bound give
\begin{equation}\label{vol:prediction-bound}
 \Pp(\text{correct label at }q_i\mid\text{past labels})\le\rho^{-N}.
\end{equation}
For completeness, if the surviving weights are $b_0,b_1,\ldots$ with
$b_{r+1}\le q b_r$ and $q=\rho^{-N}$, their tail has mass at most
$qb_0/(1-q)$, so its posterior fraction is at most $q$.
Consequently
\[
 N\ge\frac{\log(T/\eta)}{\log(1+\ell/4)}
\]
ensures, by a union bound, that sampled posterior predictions make every one
of these $T$ mistakes with probability at least $1-\eta$.
The same estimate implies a wrong posterior majority at every designated
round whenever $\rho^{-N}<1/2$.
When $H$ and $B$ are polynomial in $k$, the displayed cold width threshold is
$2^k\operatorname{poly}(k)\log(T/\eta)$.
This quantitative threshold concerns the cold posterior. A simultaneous
width-and-scale limit also needs the separate finite-scale posterior error
bound; its dimension-dependent constants should not be suppressed into this
cold width estimate.


\section{Completing the construction and the width statements}

\subsection{An explicit precision choice}

Set $p=20k+200$ and $L=21k+202$. There are $k+H$ base inputs and one
selector, so $d=k+H+1=167k+1617$. The constants in the volume argument
satisfy
\[
 C_1\le7H^2B^2,\qquad T_0\le14HB^2,\qquad C_2\le784H^4B^4.
\]
Consequently its logarithmic-slope bound satisfies
\begin{equation}\label{eq:slopeexplicit}
 \ell\ge\frac1{1568H^4B^4\,2^k}
 >2^{-54}L^{-12}2^{-k}.
\end{equation}
The required inequalities $\epsilon<1/8$ and $\epsilon<\ell/32$ hold:
the latter follows from $p>2k+63+12\log_2 L$.
Indeed $L\le223k$, so $\log_2L<8+\log_2k\le k+7$,
and $p=20k+200>14k+147$.

\subsection{The target and the fixed permutation}

Use the canonical main/helper coefficients from \eqref{eq:canonical}
and the branch bias $u_*=1/4$ on selector $z=-1$. Put
$H_\Sigma=\sum_lh_l$ and $a=2^p$. On the branch $z=+1$ choose
\[
 v_*=-H_\Sigma-R_*+a.
\]
The two branch biases are implemented by the ordinary constant and
selector coefficients $b=(u_*+v_*)/2$ and $c=(v_*-u_*)/2$.
On the all-positive helper face, the all-positive main input has score
$a>0$ and every other main input has score at most $-a<0$.
Every off-face input is negative by the helper bound. Thus the full
target is $f_0$ on the first branch and the AND of all main/helper bits
on the second branch. It is exactly the last finite cutoff function
in the volume construction.

Take the fixed input permutation to be: all $z=-1$ inputs in
lexicographic order; then the two face endpoints $q_0,q_M$ where
$M=2^k-1$; then $q_1,\ldots,q_{M-1}$; then all remaining inputs in
lexicographic order. These blocks are disjoint. The first endpoint is
negative, the last positive, and every designated middle input is
negative. After $q_1,\ldots,q_j$ have been observed, the consistent
threshold functions are exactly the cutoff functions indexed
$j,j+1,\ldots,M-1$. The first has strictly largest cone probability and
predicts $q_{j+1}$ positively. There are exactly $M-1=2^k-2$ such rounds.

Proposition~\ref{prop:gamma} proves the posterior statement of
Theorem~\ref{thm:main}. For posterior sampling, apply that statement with
per-round error allowance $\eta/(M-1)$ and use the union bound; no
independence between predictions is required.

\subsection{Teacher and Boolean circuit complexity}

The canonical full threshold representation has no zero score on the
cube. Multiply all its first-layer coefficients, including the bias,
by a sufficiently small positive scalar so that their absolute values
are at most one. The hidden sign function is unchanged. Give the single
hidden neuron output weight one. This is the claimed width-one,
weight-one teacher, with global output margin one.

For completeness the target also has a linear-size Boolean circuit;
the bounded teacher is not the only complexity guarantee. Count the
$H=O(L)$ helper bits using carry-save addition. Grouping three bits of
one weight into a full adder reduces the number of live bits by one,
so $O(H)$ constant-size full adders suffice, followed by an
$O(\log H)$ ripple addition. The common coefficient
$H_*=2^{L+2}-3$ multiplies that count by a shift and two subtractions,
using $O(L+\log H)$ gates. After conversion of signs to $0/1$ bits,
the remaining dyadic terms contain at most a constant number of source
bits per level: at most eight from the four helper-pair registers and
the main bits that replace one of those pairs. Sum positive and negative
terms separately with carry-save addition and subtract the results.
This requires $O(L)$ gates and $O(L)$-bit arithmetic. Include the fixed
bias and compare the sign. This computes $f_0$ in $O(L)$ gates. A linear
size AND and a selector multiplexer compute the other branch and the
full target. Conversely every main/helper input is essential through
the AND branch. A bounded-fan-in circuit depending on $d-1$ distinct
inputs needs $\Omega(d)$ gates. Since $L=\Theta(d)=\Theta(k)$, the circuit
size is $\Theta(d)$.

\subsection{What the width bound means}

The quantitative adjacent cone ratio from the volume argument gives a
cold-posterior threshold of $\operatorname{poly}(k)2^k$ for any fixed
confidence. A logarithmic confidence factor suffices. This means that
at every width above that threshold, taking the scalar scale $\sigma$
sufficiently small produces the designated mistakes.

Under a prescribed $\gamma>1$ and $\tau>0$, one must additionally make
\[
 N\exp[-c_dN^{\gamma-1}/(8\tau^2)]
\]
sufficiently small relative to the finite cone probabilities, as well
as suppress the nonthreshold mass in \eqref{eq:massbounds}.
These are explicit sufficient inequalities, and hold eventually for
every fixed $\gamma>1$. Since $c_d$ is of order $2^{-d}$, the resulting
fixed-$\gamma$ width threshold need not be
$\operatorname{poly}(k)2^k$. No lower bound at every width $N\ge1$ or at
every width above the width-one teacher is claimed.

\section{Polynomial IID sample complexity for the same prior}

The lower bound is specific to adversarial ordering. The same actual
posterior has a polynomial realizable IID guarantee for this target.

\begin{theorem}\label{thm:iid}
Fix any distribution $D$ on $X$, a threshold target $f$, and numbers
$\varepsilon,\beta\in(0,1)$. If
\begin{equation}\label{eq:samples}
 m\ge\frac4\varepsilon\bigl[(d+1)^2\log2+\log(1/\beta)\bigr],
\end{equation}
then, with probability at least $1-\beta$ over $m$ IID labeled examples,
posterior-label majority under \eqref{eq:gamma} has $D$-error at most
$\varepsilon$ at every sufficiently large width, for any fixed
$\gamma>1$ and $\tau>0$. The width threshold can be chosen independently
of the training sample and $D$. The assertion also holds at $\gamma=1$
for every sufficiently small fixed $\tau>0$.
\end{theorem}
\begin{proof}
The $T=2^d$ input points give $T$ hyperplanes in the $(d+1)$-dimensional
coefficient space. The usual hyperplane-region recurrence, obtained by
adding one hyperplane at a time, bounds their number of cells by
$\sum_{j=0}^{d+1}\binom Tj\le(T+1)^{d+1}\le(2T)^{d+1}$.
Every strict threshold truth table is one such cell. Hence
$\log|\cH_d|\le(d+1)^2\log2$.
A union bound shows that the probability any consistent threshold
function has $D$-error greater than $\varepsilon/4$ is at most
$|\cH_d|e^{-m\varepsilon/4}\le\beta$.

The target is always consistent. Lemma~\ref{lem:posterior} bounds the
posterior mass of all nonthreshold functions, uniformly over every
target-consistent history, by
\begin{equation}\label{eq:nonthreshold}
 B_N=\min\left\{1,\ 2^T\left(\frac{\delta_N}{p_f}\right)^N\right\}.
\end{equation}
For fixed $\gamma>1$, $B_N\to0$. For $\gamma=1$, choosing
$\delta_\tau<p_f$ also gives $B_N\to0$.
On the high-probability sample event, the posterior mean of population
error is at most $\varepsilon/4+B_N$. The population error of posterior
label majority is at most twice that mean, because a wrong majority
puts at least half the posterior mass on the wrong label. Once
$B_N\le\varepsilon/4$, the asserted bound follows.
\end{proof}

For the explicit target, $d=\Theta(k)$, so the IID sample bound is
$O((k^2+\log(1/\beta))/\varepsilon)$, while the adversarial order forces
$2^k-2$ mistakes. One sufficiently small fixed prefactor at $\gamma=1$
can satisfy both the online and IID conclusions by taking the minimum
of their positive bounds. The theorem concerns sample complexity;
the necessary learner width and posterior computation are not claimed
to be polynomial.

\section{Scope}

All labels used above are labels of one explicit affine threshold
function on the full input cube. The calibration does not pin parameters
inside a surviving truth-table cell: it eliminates incompatible cells,
and the remaining cells have different Gaussian solid angles. The
posterior comparison includes every Boolean output obtainable by the
wide sign network, including outputs that are not thresholds.
No unfavorable choice from an isotropic tie is made.

The activation is discontinuous sign. The result therefore does not
transfer to standard tanh by approximation: an exponentially narrow
threshold cone and an unbounded width limit would require new uniform
control. The exactly isotropic all-width lower-bound questions for
ordinary tanh, ReLU, and GELU remain outside this theorem.

\begin{thebibliography}{9}
\bibitem{anisotropic}
ChatGPT-6 Astra, Kaarel H\"anni, and Dmitry Vaintrob.
\emph{Exponential online mistakes at every width after arbitrarily small
Gaussian anisotropy}. 2026.
\href{../anisotropic-allwidth/paper.pdf}{Companion manuscript}.
\end{thebibliography}
\end{document}
